\section*{الفصل \ref{ch:g6:earth-a-planet} --- الأرض، كوكب من المجموعة الشمسية}

\begin{solution}{exo:g6:earth-a-planet:1}
\underline{عطارد}، \underline{الزهرة}، \underline{الأرض}،
\underline{المريخ}، المشتري، زحل، أورانوس، نبتون. \omterm{def:g6:earth-a-planet:planet-types}{والكواكب الصخرية} هي
الأربعة الأولى.
\end{solution}

\begin{solution}{exo:g6:earth-a-planet:2}
«\omterm{def:g6:earth-a-planet:planet-types}{الكوكب الصخري} صغير وكثيف، وله سطح صلب. \omterm{def:g6:earth-a-planet:planet-types}{والكوكب العملاق}
كبير جدًا، قليل الكثافة، مصنوع في معظمه من الغازات والسوائل.»
\end{solution}

\begin{solution}{exo:g6:earth-a-planet:3}
الزهرة (12\,104~كيلومترًا، مقابل 12\,756~كيلومترًا للأرض). وأكبرها
المشتري (142\,984~كيلومترًا)؛ وأصغرها عطارد (4\,879~كيلومترًا).
\end{solution}

\begin{solution}{exo:g6:earth-a-planet:4}
(أ) خطأ: لا ماء سائل على الزهرة؛ فهي حارة جدًا. (ب) صواب
(\babelsublr{$-65$~$^\circ$C} في المتوسط، مقابل \babelsublr{15~$^\circ$C}). (ج) خطأ: الكواكب
الأربعة الأقرب إلى الشمس هي \omterm{def:g6:earth-a-planet:planet-types}{الكواكب الصخرية}؛ أما \omterm{def:g6:earth-a-planet:planet-types}{الكواكب العملاقة} فهي
الأربعة الأبعد. (د) صواب: نحو 0.7~غرام لكل سنتيمتر مكعب.
\end{solution}

\begin{solution}{exo:g6:earth-a-planet:5}
$4879 \div 12756 \approx 0.382$، أي 38~\%: فعرض عطارد أكثر بقليل
من ثلث عرض الأرض.
\end{solution}

\begin{solution}{exo:g6:earth-a-planet:6}
$120536 \div 12756 \approx 9.4$: أي نحو 9 أرضين. وزحل \omterm{def:g6:earth-a-planet:planet-types}{كوكب عملاق}.
\end{solution}

\begin{solution}{exo:g6:earth-a-planet:7}
الكوكب~A \omterm{def:g6:earth-a-planet:planet-types}{كوكب عملاق}: فهو كبير جدًا، وكثافته، 1.1~غرام لكل
سنتيمتر مكعب، قليلة، مثل كثافة \omterm{def:g6:earth-a-planet:planet-types}{الكواكب العملاقة} في المجموعة الشمسية
(بين 0.7 و1.7). والكوكب~B \omterm{def:g6:earth-a-planet:planet-types}{كوكب صخري}: فهو صغير،
وكثافته، 5.0، عالية، مثل كثافة \omterm{def:g6:earth-a-planet:planet-types}{الكواكب الصخرية} (بين 3.9 و5.5).
\end{solution}

\begin{solution}{exo:g6:earth-a-planet:8}
(أ) أحر بمقدار $464 - 15 = 449$~درجة. (ب) أبرد بمقدار $15 - (-65) = 80$~درجة.
\end{solution}

\begin{solution}{exo:g6:earth-a-planet:9}
للزهرة غلاف جوي سميك جدًا من ثاني أكسيد الكربون، يحبس حرارة
الشمس كبطانية ثقيلة ولا يدعها تفلت. أما عطارد، وهو أقرب إلى
الشمس، فيكاد لا هواء له، فلا يستطيع أن يحتفظ بحرارته.
\end{solution}

\begin{solution}{exo:g6:earth-a-planet:10}
أبعد بمقدار $228 - 150 = 78$~مليون كيلومتر. $78 \div 150 = 0.52$: فبعد المريخ
أكبر من بعد الأرض بنسبة 52~\%.
\end{solution}

\begin{solution}{exo:g6:earth-a-planet:11}
قبل مليارات السنين، كانت في المريخ \omterm{def:g1:landscapes-around-us:water}{أنهار} \omterm{def:g1:landscapes-around-us:water}{وبحيرات} \omterm{def:g5:water-shapes-the-land:flood}{وفيضانات}: فقد كان أدفأ
وأكثر رطوبة. واليوم ماؤه متجمد، في جليد القلنسوتين القطبيتين وتحت
الأرض قرب \omterm{def:g3:continents-and-oceans:poles}{القطبين}. ولا يستطيع الماء السائل أن يدوم على سطحه الآن
لأن المريخ بارد جدًا (\babelsublr{$-65$~$^\circ$C} في المتوسط) ولأن هواءه
رقيق إلى حد أن الماء سرعان ما يتجمد أو يتحول إلى بخار.
\end{solution}

\begin{solution}{exo:g6:earth-a-planet:12}
(أ) $8.9 \div 9.8 \approx 0.908$، أي 91~\%. (ب) تشبه التوأم: تكاد تكون
بالحجم نفسه (95~\% من قطر الأرض)، وشدة جاذبيتها تكاد تكون نفسها،
وهي \omterm{def:g6:earth-a-planet:planet-types}{كوكب صخري} أيضًا. ولا تشبه التوأم: فهي أحر بكثير
(\babelsublr{464~$^\circ$C} مقابل \babelsublr{15~$^\circ$C})، وهواؤها أسمك بكثير ومصنوع من
ثاني أكسيد الكربون فيه \omterm{def:g4:clouds-rain-snow:cloud}{سحب} من الحمض، ولا ماء سائل فيها.
\end{solution}

\begin{solution}{pb:g6:earth-a-planet:1}
\textbf{1.} الزهرة (108~مليون~كيلومتر)، ثم الأرض (150)، ثم المريخ (228).

\textbf{2.} $150 - 108 = 42$: الزهرة أقرب بمقدار 42~مليون كيلومتر.
$228 - 150 = 78$: والمريخ أبعد بمقدار 78~مليون كيلومتر.

\textbf{3.} $6792 \div 12756 \approx 0.532$، أي 53~\%: فعرض المريخ نحو
نصف عرض الأرض.

\textbf{4.} الثلاثة كلها \omterm{def:g6:earth-a-planet:planet-types}{كواكب صخرية}: فكثافاتها (5.2 و5.5 و3.9
غرامات لكل سنتيمتر مكعب) عالية، وأقطارها صغيرة
مقارنةً بأقطار \omterm{def:g6:earth-a-planet:planet-types}{الكواكب العملاقة}.

\textbf{5.} أسمك هواء على الزهرة (92~بارًا)؛ وأرقّه على المريخ
(0.01~بار).

\textbf{6.} $1 \div 0.01 = 100$: هواء المريخ يضغط أقل من
هواء الأرض مئة مرة.

\textbf{7.} هواؤها السميك من ثاني أكسيد الكربون يحبس حرارة الشمس ولا
يدعها تفلت، كبطانية ثقيلة.

\textbf{8.} $92 \div 1 = 92$: على أرض الزهرة، يضغط الهواء
\textbf{أشد 92 مرة} مما يضغط على أرض الأرض. فيجب أن تكون للمسبار
جدران متينة جدًا لكي لا يُسحق (ويجب أيضًا أن يقاوم الحرارة
\omterm{def:g4:clouds-rain-snow:cloud}{وسحب} الحمض).

\textbf{9.} متوسط درجة حرارة الأرض، \babelsublr{15~$^\circ$C}، يقع بين
\babelsublr{0~$^\circ$C} و\babelsublr{100~$^\circ$C}: فالماء ليس متجمدًا ولا يغلي، فيبقى
سائلًا.

\textbf{10.} على الزهرة، عند \babelsublr{464~$^\circ$C}، أي فوق \babelsublr{100~$^\circ$C} بكثير، يغلي الماء
ويتبخر فورًا. وعلى المريخ، عند \babelsublr{$-65$~$^\circ$C}،
أي تحت \babelsublr{0~$^\circ$C} بكثير، يتجمد جليدًا (ويتحول ببطء إلى
بخار، لأن الهواء رقيق جدًا).

\textbf{11.} قديمًا، كان المريخ أدفأ وأكثر رطوبة من اليوم: فكان الماء السائل
يستطيع أن يجري على سطحه. وقد تغير مناخه كثيرًا: فهو الآن
\omterm{def:g5:weather-and-climate:desert}{صحراء} باردة جافة.

\textbf{12.} مثلًا: «أعزائي المسابير، هنا على الأرض عندنا ماء سائل،
في المحيطات \omterm{def:g1:landscapes-around-us:water}{والأنهار} \omterm{def:g1:landscapes-around-us:water}{والبحيرات}. ودرجات حرارتنا معتدلة، \babelsublr{15~$^\circ$C}
في المتوسط، بين تجمد الماء وغليانه. ونحن على
البعد المناسب من الشمس، لا قريبون جدًا ولا بعيدون جدًا. وغلافنا
الجوي يبقينا دافئين، ويعطي الكائنات الحية الغازات التي تحتاج إليها، ويدع
الماء يبقى سائلًا. ولهذا كانت الحياة ممكنة هنا.»
\end{solution}
